\section{Introduction to Neural Network}

\subsection*{Computation of a Single Neuron: $f(w^\top x + b)$}

A neuron computes $f(w^\top x + b)$, because the input is usually a vector.

\begin{itemize}
    \item \textbf{Single input (scalar):} $a = f(wx + b)$
    \item \textbf{Multiple inputs (vector $x \in \mathbb{R}^d$):} $z = w^\top x + b$, then $a = f(z)$
\end{itemize}

The neuron does two things: an affine step (weighted sum plus bias) that gives the pre-activation $z$, then a nonlinearity $f$ applied to that scalar. The nonlinearity is the only place a curve can enter. If $f$ were the identity, stacking layers would collapse into one affine map.

For a whole layer, the same idea is written with a matrix: $z = W a^{(l-1)} + b$, $a = f(z)$, with $f$ applied elementwise.

\subsection*{Number of Inputs to a Neuron}

A neuron can take multiple values as input. A single input is just the special case $d = 1$.

With $d$ inputs, the neuron has \textbf{one weight per input} but only \textbf{one bias}: $x \in \mathbb{R}^d$, $w \in \mathbb{R}^d$, $b \in \mathbb{R}$. It multiplies each input by its own weight, sums the results, adds the bias, and applies $f$:
\[
z = w_1 x_1 + w_2 x_2 + \dots + w_d x_d + b = w^\top x + b, \qquad a = f(z)
\]

Whatever the number of inputs, the output is always \textbf{one scalar}.

For example, in a $784 \to 128 \to 64 \to 10$ network, each neuron in the first hidden layer takes all 784 inputs, so it has 784 weights and 1 bias. This is why a layer's parameter count is $n_{\text{out}} \cdot n_{\text{in}} + n_{\text{out}}$: $n_{\text{out}}$ neurons, each with $n_{\text{in}}$ weights and one bias.

\subsection*{What Is a Layer?}

A layer is a group of neurons that all receive the same input vector and compute in parallel, each producing one scalar. Stacking those scalars gives the layer's output vector, which becomes the input to the next layer.

If a layer has $n_{\text{out}}$ neurons and receives $n_{\text{in}}$ inputs, then every neuron has its own $n_{\text{in}}$ weights and its own bias. Stacking the weight vectors as rows gives one matrix, so the whole layer computes in a single step:
\[
z = W a_{\text{prev}} + b, \qquad a = f(z)
\]
with $W$ of shape $(n_{\text{out}} \times n_{\text{in}})$, and $b$, $z$, $a$ of length $n_{\text{out}}$. The activation $f$ is applied to each coordinate independently.

\textbf{Example: a $2 \to 2 \to 1$ network.}
\begin{itemize}
    \item \textbf{Input layer:} just $x$, the raw inputs. It has no weights and does no computation.
    \item \textbf{Hidden layer:} 2 neurons, each seeing both inputs ($W^{(1)}$ is $2 \times 2$). It is called hidden because its outputs are neither the inputs nor the final prediction.
    \item \textbf{Output layer:} 1 neuron, seeing both hidden outputs ($W^{(2)}$ is $1 \times 2$), producing $\hat{y}$.
\end{itemize}

A layer where every neuron sees every input from the previous layer is called a \textbf{fully connected} (or dense) layer. A network with several of these in a row is an MLP.

\subsection*{Do Input Layer Neurons Perform Computation?}

No. The input layer's ``neurons'' are just placeholders that hold the raw input values ($x_1, x_2, \dots$). They have no weights, no bias and no activation, so nothing is computed there.

The first real computation happens in the first hidden layer, which takes the input values and applies $f(Wx + b)$. The weight matrix $W^{(1)}$ is described as connecting the input layer to the first hidden layer, but it belongs to the hidden layer's neurons: its shape is $(n_{\text{hidden}} \times n_{\text{input}})$.

This is also why a $2 \to 2 \to 1$ network is usually called a ``2-layer network'' (hidden + output). The input layer is drawn as circles for visualization, but it is not counted because it does no computation.

\subsection*{Input Layer Size for a Tabular Dataset}

Suppose a tabular dataset has 1000 rows and 4 columns. The number of input neurons depends on what the columns are.

\begin{itemize}
    \item \textbf{All 4 columns are input features (target is separate):} 4 input neurons, one per feature. The number of rows does not matter; each row is passed in as a vector $x \in \mathbb{R}^4$, and the 1000 rows are 1000 training examples.
    \item \textbf{One of the 4 columns is the target (label):} only 3 are features, so the input layer has 3 neurons. The target column is what the output layer is compared against in the loss, not something fed in.
    \item \textbf{Preprocessing can change the input size:} a categorical column with $k$ categories, one-hot encoded, becomes $k$ input neurons instead of 1. Dropping a column, or adding engineered features, changes the count too.
\end{itemize}

The rule is: \textbf{input layer size = number of numeric input values per row after preprocessing}, not the raw column count. The label is excluded.

Since there are 1000 rows, they are typically passed in batches (say 32 at a time). The input is then a matrix of shape $(\text{batch} \times n_{\text{features}})$, but the network's structure is unchanged: batch size never affects the number of neurons or weights.

\subsection*{Output Layer Design Based on the Label Type}

For a dataset with 3 numerical features and 1 label, the input layer has 3 neurons. The label is not fed in; it only appears in the loss, where it is compared against the network's output.

The output layer depends on the label type:
\begin{itemize}
    \item \textbf{Regression (continuous label):} 1 output neuron, usually with no activation (or identity).
    \item \textbf{Binary classification:} 1 output neuron with sigmoid.
    \item \textbf{Multi-class classification with $k$ classes:} $k$ output neurons with softmax.
\end{itemize}

A network such as $3 \to 8 \to 4 \to 1$ would work for regression or binary classification. The hidden sizes (8 and 4) are a free choice; only the input size (3) and the output size (set by the task) are fixed by the data.

\subsection*{One-Hot Encoding of Categorical Features}

A neuron computes $w^\top x + b$, so every input has to be a number, and the number's size has to mean something. A categorical column like color $= \{\text{red}, \text{green}, \text{blue}\}$ is just labels. If we encode them as red $= 1$, green $= 2$, blue $= 3$, we tell the network that blue is ``three times red'' and green sits ``between'' them. The weight multiplies that number, so the network is forced to treat the categories as points on a line with a fixed order and spacing, which is usually false.

\textbf{One-hot encoding} avoids this by giving each category its own 0/1 column:

\begin{center}
\begin{tabular}{lccc}
\toprule
color & is\_red & is\_green & is\_blue \\
\midrule
red   & 1 & 0 & 0 \\
green & 0 & 1 & 0 \\
blue  & 0 & 0 & 1 \\
\bottomrule
\end{tabular}
\end{center}

One original column becomes $k$ columns (here $k = 3$), and exactly one of them is 1 in each row. Since the input layer has one neuron per input value, those $k$ columns need $k$ input neurons.

This also fixes the weight problem. The first hidden neuron now has a separate weight for each category: $w_{\text{red}}, w_{\text{green}}, w_{\text{blue}}$. Its pre-activation is
\[
z = w_{\text{red}} \cdot \text{is\_red} + w_{\text{green}} \cdot \text{is\_green} + w_{\text{blue}} \cdot \text{is\_blue} + (\text{other terms}) + b
\]
For a red row, only $w_{\text{red}}$ contributes. For a green row, only $w_{\text{green}}$ does. So the network learns an independent effect for each category, with no imposed ordering or spacing.

\textbf{Example:} a dataset has 3 numerical features plus one categorical column with 4 categories. After one-hot encoding, each row is a vector of length $3 + 4 = 7$, so the input layer has 7 neurons.

\textbf{Notes:}
\begin{itemize}
    \item If a column has thousands of categories (like a zip code), one-hot becomes very wide. In that case, \textbf{embeddings} are used, which map each category to a short learned vector instead.
    \item If the categories have a natural order (small $<$ medium $<$ large), encoding them as 1, 2, 3 can be reasonable, because the ordering is real.
\end{itemize}

\subsection*{Why Each Feature Gets One Input Slot Instead of Feeding All Features into One Neuron}

The input layer is not made of computing neurons. It is just the input vector, with one slot per value, so ``one neuron per feature'' simply means ``one slot per number.'' A neuron can take multiple inputs, but that happens in the next layer, not in the input layer.

\textbf{Feeding all features into one neuron is possible, but that neuron is a hidden-layer neuron}, and it has two consequences:

\begin{enumerate}
    \item \textbf{It compresses the data.} With 3 features, one neuron computes $z = w_1 x_1 + w_2 x_2 + w_3 x_3 + b$ and outputs a single scalar. Everything the next layer sees is that one number, so many different feature combinations collapse to the same value and information is lost. For example, with $w = (1, 1, 1)$, the rows $(1, 2, 3)$ and $(3, 2, 1)$ are indistinguishable.
    \item \textbf{It limits what the network can learn.} The network can only use one weighted combination of the features. A layer with several neurons learns several different combinations (each with its own weight vector), and later layers can combine those into more complex patterns. One neuron gives only a single ``view'' of the data.
\end{enumerate}

So the structure is:
\begin{itemize}
    \item \textbf{Input layer:} holds the raw values, 3 slots, no computation, nothing lost.
    \item \textbf{First hidden layer:} every neuron takes all 3 features as input, and the number of neurons (a free choice) sets how many combinations are learned.
\end{itemize}

The input layer must keep the features separate so that the hidden neurons can each apply their own weights to each individual feature.

\subsection*{Weight Matrix Shapes: Example $4 \to 4 \to 3 \to 1$}

Consider a dataset with 1000 rows and 5 columns (4 numeric features, 1 label), and a network with architecture $4 \to 4 \to 3 \to 1$.

\textbf{Hidden layer 1.} The weight matrix has shape $(4 \times 4)$ using the convention $W$ of shape $(n_{\text{out}} \times n_{\text{in}})$. The layer has 4 neurons ($n_{\text{out}} = 4$) and receives the 4 input features ($n_{\text{in}} = 4$). Each row of $W$ holds one neuron's 4 weights, and the bias $b$ has length 4. The layer computes $z = Wx + b$, with $W$ ($4 \times 4$) times $x$ ($4 \times 1$) giving $z$ ($4 \times 1$).

\textbf{Parameters for the whole architecture:}
\begin{itemize}
    \item Hidden layer 1: $W^{(1)}$ is $4 \times 4$, $b^{(1)}$ has length 4 $\to$ $16 + 4 = 20$ parameters
    \item Hidden layer 2: $W^{(2)}$ is $3 \times 4$, $b^{(2)}$ has length 3 $\to$ $12 + 3 = 15$ parameters
    \item Output layer: $W^{(3)}$ is $1 \times 3$, $b^{(3)}$ has length 1 $\to$ $3 + 1 = 4$ parameters
\end{itemize}
Total: 39 parameters.

\subsection*{Weight Matrix Shape Convention: $(n_{\text{out}} \times n_{\text{in}})$ vs.\ $(n_{\text{in}} \times n_{\text{out}})$}

Why is $W^{(2)}$ written as $3 \times 4$ and not $4 \times 3$? It depends on the convention. Here we use $(n_{\text{out}} \times n_{\text{in}})$, where each row holds one neuron's weights.

Hidden layer 2 has 3 neurons and receives 4 inputs:
\begin{itemize}
    \item Each neuron needs 4 weights (one per input), so each neuron is a row of length 4.
    \item There are 3 neurons, so there are 3 rows.
    \item $W^{(2)}$ is therefore $3 \times 4$.
\end{itemize}

This makes the math work with the column-vector form $z = Wa + b$: $(3 \times 4)$ times $(4 \times 1)$ gives $(3 \times 1)$, a vector with one value per neuron.

The alternative $4 \times 3$ is the other common convention, $(n_{\text{in}} \times n_{\text{out}})$, where each \emph{column} holds one neuron's weights. It is valid, but the computation is then written as $z = a^\top W + b$, i.e.\ $(1 \times 4)$ times $(4 \times 3)$ gives $(1 \times 3)$. PyTorch's \texttt{nn.Linear} stores the weight as $(n_{\text{out}} \times n_{\text{in}})$, while Keras's \texttt{Dense} stores it as $(n_{\text{in}} \times n_{\text{out}})$. The two are transposes of each other, with the same 12 numbers.

Whichever convention is chosen, it must be used consistently so that the shapes multiply correctly. In general, a layer's weight shape is (neurons in this layer $\times$ neurons in the previous layer) under the first convention.